\chapter{Lymphedema}

\section*{Definition}
Lymphedema is chronic, progressive swelling of a body part due to impaired lymphatic drainage. It most commonly affects the arms or legs but can involve other areas. Primary lymphedema is congenital or genetic; secondary lymphedema results from damage to the lymphatic system.

\section*{Pathophysiology}
The lymphatic system returns interstitial fluid to the bloodstream. When lymph vessels or nodes are damaged or absent, protein-rich fluid accumulates in the interstitial space. This causes swelling, tissue fibrosis, chronic inflammation, and impaired immune function. The accumulated fluid provides a growth medium for bacteria, increasing infection risk.

\section*{Causes}
\begin{itemize}
\item Secondary (most common): cancer treatment (lymph node dissection, radiation)
\item Breast cancer-related lymphedema (arm swelling after axillary dissection)
\item Pelvic cancer treatment (leg swelling)
\item Infection: filariasis (tropical, caused by Wuchereria bancrofti)
\item Infection: recurrent cellulitis or lymphangitis
\item Trauma or burns
\item Venous insufficiency (chronic, with secondary lymphatic overload)
\item Primary (idiopathic): congenital or pubertal onset (Milroy disease, Meige disease)
\item Obesity-associated lymphedema
\end{itemize}

\section*{When to See a Doctor}
\begin{itemize}
\item Persistent swelling of an arm or leg
\item Swelling that begins distally and progresses proximally
\item Non-pitting edema in later stages (fibrosis)
\item Heaviness, tightness, or fullness in the affected limb
\item Decreased flexibility in the affected joint
\item Aching or discomfort
\item Skin changes: thickening, hyperkeratosis, papillomatosis
\item Positive Stemmer sign (inability to pinch the skin on the dorsum of the second toe)
\end{itemize}

\section*{Related Symptoms}
\begin{itemize}
\item Edema (generalized fluid retention)
\item Venous insufficiency
\item Cellulitis (recurrent infections)
\item Lymphangitis
\item Elephantiasis (severe filariasis)
\end{itemize}

\section*{Self-Care/Home Management}
\begin{itemize}
\item Elevate the affected limb when possible
\item Practice good skin hygiene and moisturize daily
\item Protect the limb from cuts, burns, and insect bites
\item Wear compression garments as prescribed
\item Perform light exercise with the affected limb
\item Avoid constrictive clothing or jewelry on the affected side
\item Avoid prolonged standing or sitting
\item Maintain a healthy body weight
\end{itemize}

\section*{How Your Doctor Will Evaluate You}
\begin{itemize}
\item Clinical examination: pitting, Stemmer sign, skin changes
\item Lymphoscintigraphy (nuclear medicine imaging of lymphatic flow)
\item Bioimpedance spectroscopy for early detection
\item MRI or CT for detailed anatomy
\item Duplex ultrasound to rule out venous causes
\end{itemize}

\section*{Treatment Options}
\begin{itemize}
\item Complete decongestive therapy: manual lymphatic drainage, compression bandaging, exercise, skin care
\item Compression garments for maintenance
\item Pneumatic compression devices
\item Surgical options: lymphovenous anastomosis, vascularized lymph node transfer
\item Liposuction for late-stage fibrotic lymphedema
\item Antibiotics for cellulitis prophylaxis
\item Avoid blood pressure measurements and venipuncture on the affected arm
\end{itemize}
